% =====================================================================
% Приложение Г. Открытый исходный код (четвертое приложение работы,
% буквенная нумерация А, Б, В, Г из main.tex).
% Ссылка на репозиторий + согласие на публикацию ВКР
% (требование методички, раздел 5.5 и 5.6.3).
% =====================================================================

\clearpage
\section{Открытый исходный код}
\label{app:github}

В соответствии с пунктом 5.5 Правил подготовки и защиты ВКР
полная программная реализация настоящего исследования размещена
в открытом доступе на платформе GitHub.

\begin{description}[leftmargin=2em, itemsep=0pt]
\item[\textbf{Адрес репозитория}]
   \url{https://github.com/cactus-tim/adaptive-topologies-mas}
\item[\textbf{Идентификатор фиксации}] коммит
   \code{058271a39a0e} (полный хеш ---
   \mbox{058271a39a0e4ecc0e3b47bd1d826c03fc33872b}) ветки
   \code{main}.
\item[\textbf{Лицензия}] MIT (файл \code{LICENSE}).
\item[\textbf{Сборка и воспроизведение}] инструкция в
   \code{README.md}; конфигурации экспериментов главы~\ref{sec:experiments}
   --- в \code{conf/experiments/}, запуск через
   \code{atm run~/~atm grid --config conf/experiments/e1.yaml}.
\end{description}

Согласие автора на публикацию ВКР на портале НИУ ВШЭ (п. 5.6.3
Правил) оформлено QR-кодом системы «LMS-Антиплагиат» в составе
пакета документов в ГЭК.
